\documentclass[a4paper]{article}

\usepackage{graphicx}
\usepackage{amsmath,amssymb}
\usepackage{minted}
\usepackage{multirow}
\usepackage{float}
\usepackage{todonotes}
\usepackage{forest}
\usepackage{xstring}
\usepackage{xcolor}
\usepackage[most]{tcolorbox}
\usepackage{xparse}
\usepackage{natbib}

\usetikzlibrary{graphs,graphdrawing}
\usegdlibrary{layered}

% for external graphviz
\input{/home/donkarlo/Dropbox/repo/nd_ai_project/src/nd_ai/knoweldge_representation/language/formal/computer/programming/tex/latex/code_clips/external_graphviz.tex}

% for tags
\input{/home/donkarlo/Dropbox/repo/nd_ai_project/src/nd_ai/knoweldge_representation/language/formal/computer/programming/tex/latex/parts/control_sequence/kind/semtag/semtag.tex}

% for links
\input{/home/donkarlo/Dropbox/repo/nd_ai_project/src/nd_ai/knoweldge_representation/language/formal/computer/programming/tex/latex/code_clips/seehere.tex}

\title{Neural summation}
\date{}
\author{Mohammad Rahmani}

\begin{document}
\maketitle
    \seeherefooter{https://en.wikipedia.org/wiki/Summation_(neurophysiology)}
    \seeherefooter{https://www.lecerveau.ca/flash/capsules/articles_pdf/temporal_summation.pdf} defines neural summation as Summation, which includes both spatial summation and temporal summation, is the process that determines whether or not an action potential will be generated by the combined effects of excitatory and inhibitory signals, both from multiple simultaneous inputs (spatial summation), and from repeated inputs (temporal summation). Depending on the sum total of many individual inputs, summation may or may not reach the threshold voltage to trigger an action potential.
    \bibliographystyle{plainnat}
    \bibliography{/home/donkarlo/Dropbox/repo/nd_shared_research/bibliography/bibliography.bib}
\end{document}